%File: ext-abstract.tex
\documentclass[letterpaper]{article} % DO NOT CHANGE THIS
\usepackage[submission]{aaai2026}  % DO NOT CHANGE THIS
\usepackage{times}  % DO NOT CHANGE THIS
\usepackage{helvet}  % DO NOT CHANGE THIS
\usepackage{courier}  % DO NOT CHANGE THIS
\usepackage[hyphens]{url}  % DO NOT CHANGE THIS
\usepackage{graphicx} % DO NOT CHANGE THIS
\urlstyle{rm} % DO NOT CHANGE THIS
\def\UrlFont{\rm}  % DO NOT CHANGE THIS
\usepackage{natbib}  % DO NOT CHANGE THIS AND DO NOT ADD ANY OPTIONS TO IT
\usepackage{caption} % DO NOT CHANGE THIS AND DO NOT ADD ANY OPTIONS TO IT
\frenchspacing  % DO NOT CHANGE THIS
\setlength{\pdfpagewidth}{8.5in} % DO NOT CHANGE THIS
\setlength{\pdfpageheight}{11in} % DO NOT CHANGE THIS

\usepackage{algorithm}
\usepackage{algorithmic}
\usepackage{verbatim}
\usepackage{tikz}
\usepackage{amssymb,amsmath}

\usepackage{newfloat}
\usepackage{listings}
\DeclareCaptionStyle{ruled}{labelfont=normalfont,labelsep=colon,strut=off} % DO NOT CHANGE THIS
\lstset{%
	basicstyle={\footnotesize\ttfamily},%
	numbers=left,numberstyle=\footnotesize,xleftmargin=2em,%
	aboveskip=0pt,belowskip=0pt,%
	showstringspaces=false,tabsize=2,breaklines=true}
\floatstyle{ruled}
\newfloat{listing}{tb}{lst}{}
\floatname{listing}{Listing}

\pdfinfo{
/TemplateVersion (2026.1)
}

\setcounter{secnumdepth}{0}

\title{Agentic Engineering for Quantum Algorithm Design (Extended Abstract)}
\author{Anonymous Submission}
\affiliations{}

\begin{document}

\maketitle

\begin{abstract}
Large language models (LLMs) have shown competence in abstract mathematical reasoning.
Some frontier models have even been able to solve open problems in combinatorics.
However, these models often fail to consistently produce correct, verifiable results, and require
much human intervention and iteration.
Furthermore, LLMs struggle in the niche domain of quantum computing, where the potential training material is
limited compared to a more general and well-studied field such as combinatorics.
We present a work-in-progress agentic framework for abstract quantum computing algorithm design.
Core definitions and axioms of the field are defined in Lean, as are goals,
and correctness is checked at every step with the Lean proof-checker.
\end{abstract}

\section{Introduction}

Quantum computing has the potential to exponentially speed up
many classical computing algorithms \cite{brandao2010exponential}.
However, solutions to quantum computing problems are often hard to verify,
given that quantum programs cannot be efficiently simulated on a classical computer,
and access to currently-available quantum machines is very expensive.
Because of this, correctness of quantum algorithms can often only be proved analytically.
For these reasons, quantum algorithm design presents a particular challenge for LLM agents.

\section{Proposed Approach}

\paragraph{Formal Verification}

\paragraph{Agentic Environment}

\paragraph{Reinforcement Learning}

\section{Benchmarking}

\bibliography{refs}

\end{document}
